%==============================================================================
\section{GLP}
\label{sec:glp}
%==============================================================================

We present GLP syntax and operational semantics and discuss its programming techniques.

%------------------------------------------------------------------------------
\subsection{Syntax}
\label{sec:glp-ext}
%------------------------------------------------------------------------------

GLP extends LP by adding a paired \emph{reader} $X?$ to every ``ordinary'' logic variable $X$, now called a \emph{writer}.

\begin{definition}[GLP Variables]
\label{def:glp-variables}
Let $\calV$ denote the set of LP variables (identifiers beginning with uppercase), henceforth called \temph{writers}. Define $\calV? = \{X? \mid X \in \calV\}$, called \temph{readers}. The set of all GLP variables is $\hat\calV = \calV \cup \calV?$. A writer $X$ and its reader $X?$ form a \temph{variable pair}.
\end{definition}
GLP terms, unit goals, goals, and clauses are as in LP but defined over the variables in $\hat\calV$.

\begin{definition}[Single-Occurrence (SO) Invariant]
\label{def:so-invariant}
A term, goal, or clause satisfies the \temph{single-occurrence (SO) invariant} if every variable occurs in it at most once.
\end{definition}

\begin{definition}[GLP Program, Goals]
\label{def:glp-program}
A clause $C$ satisfies the \temph{single-reader/single-writer (SRSW) restriction} if it satisfies SO and a variable occurs in $C$ iff its paired variable also occurs in $C$.
A \temph{GLP program} is a finite sequence of clauses satisfying SRSW; clauses for the same predicate form a \temph{procedure}.
The set of GLP goals $\hat\calG(P)$ includes all goals over $\hat\calV$ and the vocabulary of $P$ that satisfy SO.
\end{definition}
The purpose of the SRSW restriction is to prevent multiple writer occurrences racing to assign a variable; the simpler (SW) restriction is insufficient, since if a reader with multiple occurrences is assigned a term with a writer, the result is a writer with multiple occurrences.
% We use set notation also when referring to multisets.
\begin{example}[Fair Merge]
\label{ex:merge}
Consider the quintessential concurrent logic program for fairly merging two streams, written in GLP:\footnote{Moded-type definitions (\texttt{T ::= ...}) and declarations (\texttt{procedure ...}) used informally in examples are formally introduced in the Typed GLP companion paper~\cite{shapiro2026types}.}
\begin{verbatim}
Stream ::= [] ;[_|Stream].

procedure merge(Stream?, Stream?, Stream).
merge([X|Xs], Ys, [X?|Zs?]) :- merge(Ys?, Xs?, Zs).
merge(Xs, [Y|Ys], [Y?|Zs?]) :- merge(Xs?, Ys?, Zs).
merge(Xs, [], Xs?).
merge([], Ys, Ys?).
\end{verbatim}
and the goal \verb=merge([1,2,3|Xs?],[a,b|Ys?],Zs)=. The goal satisfies SO and each clause satisfies SRSW. The first clause swaps inputs in the recursive call, ensuring fairness when both streams are available.
\end{example}


As we shall see (Proposition~\ref{prop:so-preservation}), the SO invariant is maintained by the SRSW restriction: reducing a goal satisfying SO with a clause satisfying SRSW results in a goal satisfying SO. 

%------------------------------------------------------------------------------
\subsection{Operational Semantics}
\label{sec:glp-operational}
%------------------------------------------------------------------------------

\begin{definition}[Transition System]
\label{def:ts}
A \temph{transition system} is a tuple $TS = (C, c_0, T)$ where $C$ is an arbitrary set of \temph{configurations}, $c_0 \in C$ is a designated \temph{initial configuration}, and $T \subseteq C \times C$ is a \temph{transition relation}, with transitions written $c \rightarrow c' \in T$.
A transition $c \rightarrow c' \in T$ is \temph{enabled} from configuration $c$. A configuration $c$ is \temph{terminal} if no transitions are enabled from $c$. A \temph{computation} is a (finite or infinite) sequence of configurations where for each two consecutive configurations $(c,c')$ in the sequence, $c \rightarrow c' \in T$. We write $c \xrightarrow{*} c'$ to denote the existence of a computation (empty if $c = c'$) from $c$ to $c'$. A \temph{run} is a computation starting from $c_0$, which is \temph{complete} if it is infinite or ends in a terminal configuration. The \temph{outcome} of a complete run is determined by a domain-specific function from complete runs to an outcome space.
\end{definition}


\begin{definition}[Substitutions and Assignments]
\label{def:writers-assignment}
A GLP \temph{writer assignment} is a term of the form $X := T$, $X\in\calV$, $T\notin\calV$, satisfying SO. Similarly, a GLP \temph{reader assignment} is a term of the form $X? := T$, $X?\in\calV?$, $T\notin\calV$, satisfying SO. A \temph{writers (readers) substitution} $\sigma$ is the substitution implied by a set of writer (reader) assignments that jointly satisfy SO. Given a writers assignment $X := T$, its \temph{readers counterpart} is $X? := T$, and given a writers substitution $\sigma$, its \temph{readers counterpart} $\sigma?$ is the readers substitution defined by $X?\sigma? = X\sigma$.
\end{definition}

\begin{definition}[GLP Renaming, Renaming Apart]
\label{def:glp-renaming}
A \temph{GLP renaming} is an injective substitution $\rho: \hat\calV \to \hat\calV$ such that for each $X \in \calV$: $X\rho \in \calV$ and $X?\rho = (X\rho)?$. Two GLP terms \temph{have a variable in common} if for some writer $X \in \calV$, either $X$ or $X?$ occurs in both. A GLP renaming $\sigma$ \temph{renames $T'$ apart from} $T$ if $T'\sigma$ and $T$ have no variable in common.
\end{definition}

\begin{definition}[Writer MGU]
\label{def:writer-mgu}
Given two GLP unit goals $A$ and $H$, a \temph{writer mgu} is a writers substitution $\sigma$ such that $A\sigma = H\sigma$ and $\sigma$ is most general among such substitutions. 
\end{definition}

\begin{definition}[GLP Goal/Clause Reduction]
\label{def:glp-reduction}
Given GLP unit goal $A$ and clause $C$, with $H$ \verb|:-| $B$ being the result of the GLP renaming of $C$ apart from $A$, the \temph{GLP reduction} of $A$ with $C$ \temph{succeeds with result} $(B,\sigma)$ if $A$ and $H$ have a writer mgu.
\end{definition}

\begin{definition}[cGLP Transition System]
\label{def:cglp-ts}
Given a GLP program $P$, an \temph{asynchronous resolvent} over $P$ is a pair $(G, \sigma)$ where $G \in \hat\calG(P)$ and $\sigma$ is a readers substitution.

A transition system $cGLP = (\calC, c_0, \calT)$ is a \temph{cGLP transition system} over $P$ and initial goal $G_0$ satisfying SO if:
\begin{enumerate}
    \item $\calC$ is the set of all asynchronous resolvents over $P$
    \item $c_0 = (G_0, \emptyset)$
    \item $\calT$ is the set of all transitions $(G, \sigma) \rightarrow (G', \sigma')$ satisfying either:
    \begin{enumerate}
        \item \textbf{Reduce:} there exists unit goal $A \in G$ such that $C \in P$ is the first clause for which the GLP reduction of $A$ with $C$ succeeds with result $(B, \hat\sigma)$, $G' = (G \setminus \{A\} \cup B)\hat\sigma$, and $\sigma' = \sigma \circ \hat\sigma?$
        \item \textbf{Communicate:} $\{X? := T\} \in \sigma$, $X?\in G$, $G' = G\{X? := T\}$, and $\sigma' = \sigma$
    \end{enumerate}
\end{enumerate}
\end{definition}

cGLP Reduce differs from LP in (1) the use of a writer mgu instead of a regular mgu and (2) the choice of the first applicable clause instead of any clause. The first is the fundamental use of GLP readers for communication and synchronization. The second compromises on the or-nondeterminism of LP to allow writing fair concurrent programs, such as fair merge above. Note that or-nondeterminism is not completely eliminated, as different scheduling of arrival of assignments on the two input streams of \verb|merge| may result in different orders in its output stream.

The cGLP Communicate rule realizes the use of reader/writer pairs for asynchronous communication: it communicates an assignment from its writer to its paired reader.


\mypara{Persistence of Reducibility}
Key differences between LP and cGLP relate to persistence. In LP, if a goal cannot be reduced, it will never be reduced. In cGLP, a goal that cannot be reduced now may be reduced in the future: if $A$ and $H$ have an mgu that writes on a reader $X? \in A$, and therefore have no writer mgu at present,  another goal that has $X$ may reduce, assigning $X$, and later $X?$, to a value that will allow $A$ and $H$ to have a writer mgu. Conversely, in LP, if a goal $A$ can be reduced now with some clause $H$\verb|:-|$B$, with a regular mgu of $A$ and $H$, it may not be reducible in the future due to variables that $A$ shares with other goals being assigned values by reductions of other goals, preventing unification between the instantiated $A$ and $H$. In cGLP, if a goal $A$ can be reduced now (with a writers mgu), it can always be reduced in the future, as the SO invariant ensures that no other goal can assign any writer in $A$.

Implementation-wise, if a GLP goal $A$ cannot be reduced now, but there is a readers substitution $\sigma$ such that $A\sigma$ can be reduced, such readers are identified, the goal $A$ \emph{suspends} on these readers, and is rescheduled for another reduction attempt once any of them is assigned.

Despite these differences, cGLP adopts the same notions of proper run, successful run, and outcome as LP\ifappendix{} (Definition~\ref{def:proper-run}, Appendix~\ref{appendix:lp})\fi, and has the same notion of logical consequence as LP. Let $/?$ be an operator that replaces every reader by its paired writer.

\begin{restatable}[cGLP Computation is Deduction]{proposition}{propGLPDeduction}
\label{prop:cglp-deduction}
Let $(G_0$ \texttt{:-} $G_n)\sigma$ be the outcome of a proper cGLP run $\rho: (G_0,\sigma_0) \rightarrow \cdots \rightarrow (G_n, \sigma_n)$ of $cGLP$. Then $(G_0$ \texttt{:-} $G_n)\sigma/?$ is a logical consequence of $P/?$.
\end{restatable}

We note an additional safety property of cGLP runs.

\begin{restatable}[SO Preservation]{proposition}{propSOPreservation}
\label{prop:so-preservation}
If the initial goal $G_0$ satisfies SO, then every goal in the cGLP run satisfies SO.
\end{restatable}

Fairness requires identifying ``the same transition'' across configurations as goals are instantiated by reader substitutions. Following~\cite{lewis2026volitional}, we define goal identity and transition equivalence.

\begin{definition}[Goal Identity]
\label{def:goal-identity}
Each unit goal in a cGLP computation carries a unique \temph{identifier} assigned at spawn time: goals in the initial goal $G_0$ receive distinct initial identifiers, and goals spawned by a Reduce transition receive fresh identifiers. Communicate transitions preserve identifiers: the goal containing $X?$ retains its identity after instantiation.
\end{definition}

\begin{definition}[Transition Equivalence]
\label{def:transition-equivalence}
Given a set of transitions $T$, a \temph{transition equivalence}~\cite{lewis2026volitional} is an equivalence relation $\sim$ on $T$. We write $[t]$ for the equivalence class of $t$ under $\sim$. An equivalence class $[t]$ is \temph{enabled} in configuration $c$ if some $t' \in [t]$ is enabled in $c$.

Define the labelling function $\ell: \calT \to L$ on cGLP transitions by:
\[
\ell(t) :=
\begin{cases}
\langle \mathsf{Reduce},\, \mathit{id}(A),\, C \rangle & \text{if $t$ is a Reduce of unit goal $A$ with clause $C$,} \\[2pt]
\langle \mathsf{Communicate},\, X := T \rangle & \text{if $t$ is a Communicate of $X? := T$,}
\end{cases}
\]
where $\mathit{id}(A)$ is the identifier of $A$ (Definition~\ref{def:goal-identity}). The \temph{cGLP transition equivalence} is the relation $\sim\,\subseteq\,\calT\times\calT$ defined by $t_1 \sim t_2 \iff \ell(t_1) = \ell(t_2)$.
\end{definition}

Two equivalent Reduce transitions may operate on different instantiations of the same goal---with different reader substitutions applied---but produce the same writer mgu (since the writer mgu assigns only writers and is unaffected by reader substitutions).

\begin{definition}[cGLP Fair Run]
\label{def:cglp-fair-run}
A complete cGLP run is \temph{fair} if for every equivalence class $[t] \in T/{\sim}$: whenever $[t]$ becomes enabled, eventually some $t' \in [t]$ is taken.
\end{definition}

The SO invariant of GLP allows eschewing unification in favour of \emph{term matching}: if two terms that jointly satisfy SO are unifiable, their mgu maps each variable in one term to a subterm of the other. Term matching thus performs joint term-tree traversal and collects variable assignments along the way; the detailed definition and table appear in \appref{appendix:term-matching}{\arxivref}.

%------------------------------------------------------------------------------
\subsection{Guards}
\label{sec:guards}
%------------------------------------------------------------------------------

GLP clauses may include \emph{guards}---tests that determine clause applicability.

\begin{definition}[Guarded Clause]
\label{def:guarded-clause}
A \temph{guarded clause} has the form $H$ \verb|:-| $G$ \verb"|" $B$, where $H$ is the head, $G$ is a conjunction of guard predicates, and $B$ is the body. The guard separator ``\verb"|"'' distinguishes guards from the body, and is interpreted logically as a conjunction.  Guard arguments are readers paired to head writers.
\end{definition}

Guards have three-valued semantics. Each guard predicate explicitly defines its \emph{success} condition. A guard \emph{suspends} if it does not succeed but some instance of it under a readers substitution would succeed. A guard \emph{fails} if no such instance exists. A guard conjunction succeeds if all members succeed; it suspends if any member suspends and none fail; it fails if any member fails.

Definition~\ref{def:glp-reduction} of a GLP goal/clause reduction is augmented to succeed if the guard also succeeds.

\begin{remark}[Guards and SRSW]
\label{rem:guards-srsw}
Guard occurrences count toward SRSW satisfaction: if $X?$ occurs in a guard, its paired writer $X$ must occur in the head and $X?$ may additionally occur once in the body.
%
Furthermore, if the success of a guard implies that $X?$ is bound to a ground term, then both $X$ and $X?$ may occur multiple times in the clause. Groundness-implying guards include \verb|ground|, \verb|integer|, \verb|number|, \verb|string|, \verb|constant|, arithmetic comparisons (\verb|<|, \verb|>|, \verb|=<|, \verb|>=|, \verb|=:=|, \verb|=\=|), and ground equality (\verb|=?=|). However,  \verb|known| and \verb|compound| do not imply groundness.
\end{remark}

\begin{remark}[Anonymous Variables]
\label{rem:anonymous-variables}
An \emph{anonymous variable} is any variable whose name begins with \verb|_| (e.g., \verb|_|, \verb|_Out|). Anonymous variables may appear anywhere a writer variable may appear; each occurrence denotes a fresh writer with no paired reader. Anonymous readers (\verb|_?|, \verb|_Name?|) are not permitted.
\end{remark}

Guard predicates include type tests (\verb|integer|, \verb|number|, \verb|string|, \verb|constant|, \verb|compound|, \verb|ground|, \verb|known|), arithmetic comparisons (\verb|<|, \verb|>|, \verb|=<|, \verb|>=|, \verb|=:=|, \verb|=\=|), and ground equality (\verb|=?=|).

\mypara{System predicates and body kernels}
GLP \emph{system predicates} such as \verb|:=| (arithmetic assignment), \verb|now| (clock access), and \verb|=..| (term composition/decomposition) are predicates defined by GLP clauses whose bodies may invoke \emph{body kernel predicates}---runtime-implemented primitives not directly accessible to user programs. The system predicate's own guards ensure that body kernel preconditions are met before invocation. For example, arithmetic assignment is defined recursively by clauses that pattern-match on each operator:

\begin{verbatim}
Result? := N :- number(N?) | Result = N?.
Result? := X + Y :- number(X?), number(Y?) |
    '_add'(X?, Y?, Result).
Result? := X + Y :- otherwise |
    X1 := X?, Y1 := Y?, Result := X1? + Y1?.
\end{verbatim}
\noindent The base case binds a plain number; the first addition clause guards that both operands are numbers and invokes the \verb|'_add'| body kernel; the fallback clause recursively evaluates subexpressions. The complete definitions of all system predicates and the full table of body kernels appear in \appref{appendix:guards-system}{\arxivref}.

\subsection{GLP Programming Techniques}

On the one hand, GLP has no unification, only pattern-matching and single-writer variables. On the other hand, goal reduction may write on several variables atomically. The result is that almost the entire range of concurrent logic programming techniques developed since the 1980s~\cite{shapiro1989family,shapiro1987concurrent} is available in GLP, including: streams and their fair, biased, and self-balancing merge~\cite{shapiro1984fair,shapiro1986multiway}; channels as generalised streams~\cite{tribble1988channels}; messages with reply variables and bounded buffers~\cite{takeuchi1988bounded}; network reconfiguration and distributed programming~\cite{shafrir1988distributed}; object-oriented programming with mutable state via stream-recursion~\cite{shapiro1983object}; systems programming and computation control~\cite{shapiro1984systems,silverman1988logix}; and the full range of meta-interpreters --- fail-safe, termination-detecting, tracing, and algorithmic debugging~\cite{safra1988meta,lichtenstein1988concurrent}.

Two techniques require workarounds. The use of logic variables for mutual exclusion --- standard in concurrent logic languages with atomic unification~\cite{shapiro1989family} --- must be achieved in GLP via a monitor~\cite{hoare1974monitors,hansen1973operating}, the standard distributed-computing solution. Similarly, general broadcast --- distributing a value to multiple consumers --- is not directly supported by GLP, since it would require multiple occurrences of the same reader, in violation of SO; it must be achieved by explicit \emph{distributors}.  We note that if moded-type checking can determine that the values to be distributed never contain writers, then the SO restriction on readers may safely be relaxed allowing broadcast natively; a possible future upgrade of the language.  Excluded is the encoding of logic programs or-parallel by the and-parallelism of Concurrent Prolog~\cite{codish1986compiling,shapiro1989or}, which depends on general unification in a fundamental way.




